\section{A 120-qubit experiment: specification, prediction and criteria}\label{app:hw100}
The estimates of \cref{sec:sota} concern one core at a time. To use a chip of more than $100$ qubits we use the independence that the interface gives: a portfolio of $10$ assets, $6$ periods and $2$ sectors ($n=120$ variables, a looser family with at least $10$ feasible strings per block, \path{negsearch/hw100.py}) splits into twelve blocks of $10$ qubits, each with its own quadratic objective (the neighbours enter as linear fields), and the twelve blocks are placed on disjoint connected regions of $11$ physical qubits and run as one circuit of $120$ measured qubits and $120$ classical bits (a chip with at least $132$ usable qubits, such as the $156$-qubit \texttt{ibm\_fez}). Twelve circuits are specified: our light and full $p=1$ layers, our preparation alone, a Hadamard control, three baselines from the warm-started product state of the Fourier-LCU paper (single-basis Fourier-LCU, unbalanced penalty, XY-mixer QAOA; trained by exact simulation, with the same decomposition and count constraints, the setting most favourable to them) and the five cost-informed circuits of \cref{sec:core}. \Cref{tab:hw100_pred,tab:hw100_cost} give the compiled counts, depths and the noise-model prediction (\cref{sec:sota}). A complete shot is feasible with probability roughly the product of the block probabilities, so the useful output is the per-block post-selected samples, assembled into feasible assignments of all $120$ variables. The pre-stated criteria are: (1)~the preparation is above the upper $95\%$ bound of the control; (2)~preparation and light layer are above the upper bounds of the three baselines, and (2c)~the tilted preparation is above the three baselines that receive the same classical information; (3)~all twelve blocks return at least $50$ accepted samples for the preparation; (4)~the light layer has higher quality than the preparation alone; (5)~the measured feasibility of the preparation is within $0.1$ of the prediction. Criterion 2 is predicted to hold for the preparation (and 2c for the tilted one) and to fail for the light layer against the unbalanced circuit, criterion 4 may fail, and the comparison of the quality of the tilted preparation with the biased baselines (4b) is predicted to fail ($0.79$ against $0.83$--$0.88$); the comparison of quality with the baselines is reported, not required. What the experiment would show is feasibility by construction at a scale at which the global-constraint methods have no chance, with the baselines given the decomposition; it would not show a quantum advantage, since blocks of ten qubits are classically trivial, nor better optimisation inside the blocks. The notebooks (\path{notebooks/HW100_*.ipynb}) run the same analysis on the noise model, on a gate-level simulation of each block, or on a device.

\begin{table*}[t]\centering\scriptsize\setlength{\tabcolsep}{3pt}
\caption{\orig{compiled counts, uniform error model $e_2=0.0035$ per two-qubit gate; no noise simulation, nothing executed} One block of the $120$-qubit portfolio ($(l,s)=(2,1)$, $|F|=30$, 10 qubits) compiled to different connectivities with CZ as the only two-qubit gate. Entries: two-qubit gates / two-qubit depth / ESP, best of six transpiler seeds. Preparation: our preparation alone; layer: our light $p=1$ layer (\texttt{architectures.py}; a second block gives the same pattern).}\label{tab:arch}
\begin{tabular}{lcccc}\toprule
topology & our prep. & our layer & Fourier-LCU / unbalanced & XY mixer\\
\midrule
all-to-all & 90 / 64 / 0.66 & 215 / 118 / 0.38 & 90 / 34 / 0.68 & 106 / 36 / 0.62\\
grid $2{\times}5$ & 131 / 90 / 0.56 & 333 / 181 / 0.23 & 112 / 42 / 0.60 & 126 / 44 / 0.56\\
ring & 168 / 110 / 0.49 & 393 / 227 / 0.18 & 152 / 71 / 0.51 & 178 / 74 / 0.45\\
heavy-hex patch & 190 / 164 / 0.44 & 415 / 251 / 0.16 & 140 / 77 / 0.53 & 177 / 85 / 0.45\\
line & 196 / 164 / 0.43 & 418 / 242 / 0.16 & 152 / 75 / 0.51 & 180 / 70 / 0.44\\
\bottomrule\end{tabular}\end{table*}

\paragraph{Does the architecture matter?} \orig{compiled counts, uniform error model; no noise simulation} \Cref{tab:arch} compiles one block to five connectivities (\path{experiments/architectures.py}). The baselines' circuits are a product-state layer followed by one $ZZ$ layer, so their cost moves from $90$ two-qubit gates (all-to-all) to $112$--$152$ on sparse devices ($1.2$--$1.7\times$). Our preparation is more sensitive to connectivity: for the block shown it goes from $90$ (all-to-all) to $131$ on a $2\times5$ grid, $168$ on a ring and $190$ on a heavy-hex patch ($2.1\times$), because the controlled-swap chain moves units across the block; our light layer goes from $215$ to $333$ on the grid and $415$ on heavy-hex. So the ratio of our layer to the baselines is $2.4\times$ all-to-all and $2.7$--$3.0\times$ on heavy-hex, and higher connectivity (a grid or all-to-all) would lower our ESP penalty more than theirs: ESP of the layer $0.16$ (heavy-hex) against $0.23$ (grid) and $0.38$ (all-to-all), against $0.53$, $0.60$ and $0.68$ for the baselines. A native fractional $R_{ZZ}$ gate halves the baselines on all-to-all ($45$ against $90$) but changes ours by only $19\%$ ($175$ against $215$) and the sparse cases little, since routing dominates there; we did not model its error. Smaller preparations, with $|F|=10$, are nearly topology-independent ($23$--$29$ gates), so the sensitivity grows with the number of units in the block. These counts bound what a different architecture (or a better layout for the Dicke hub) could give; the conclusions of the paper concern the heavy-hex Heron devices.
